\documentclass{article}
\usepackage{graphicx} % Required for inserting images
\usepackage{amssymb}

\title{Chapter 24 Problem Set}
\author{Siddharth K}
\date{April 2024}

\begin{document}

\maketitle

\textbf{1a.} Let's consider a function $g: T \rightarrow S$ that has the same input values as the output values of $f$ and vice versa. Since $f$ is injective, every element in $S$ is an output value of $g$, which is the definition of surjectivity.
\newline

\textbf{1b}. By the previous problem, an injective function $f: S \rightarrow T$ implies the existence of a surjective function $g: T \rightarrow S$. Therefore, having injective functions from $S$ to $T$ and $T$ to $S$ means that there are surjective functions from $T$ to $S$ and $S$ to $T$. This is the condition for CSB, so it holds here as well.
\newline

\textbf{2.} An algebraic number is a root of a polynomial, but different polynomials can have the same root. So the cardinality of the set of algebraic numbers is less than the cardinality of the set of all polynomials' roots. That means, if we can find a bijection between that set and natural numbers, the set of algebraic numbers must also be countable.
\newline

Given a polynomial's root, we can take its degree and which root it is in increasing order. We can plot those two points on a graph. Doing this for all polynomials and their roots, we get a subset of $\mathbb{N} \times \mathbb{N}$. This is an injection because every input has an output in $\mathbb{N} \times \mathbb{N}$ where $|y| \le x$ and $x \ge 0$. This means the cardinality of the set of polynomial's roots is less than that of $\mathbb{N}\times \mathbb{N}$ where $|y| \le x$ and $x \ge 0$.
\newline

We can find an injection from that set to $\mathbb{N} \times \mathbb{N}$ just by mapping each element to itself, so its cardinality is less than or equal to that of $\mathbb{N} \times \mathbb{N}$, which is just aleph-zero. Therefore, the set of all polynomials' roots is countable, and therefore, so is the set of algebraic numbers.
\newline

\textbf{3.} First, we can establish an injection from $(0, 1)$ to the real numbers by mapping each number to itself. Then, we can map a real number, x, to $(0, 1)$ by finding $\frac{1}{e^x + 1}$. By CBT, these two injections imply the existence of a bijection.
\newline

We can establish an injection from $[0, 1]$ to the real numbers by mapping each number to itself. Then, we can map a real number, x, to $[0, 1]$ by finding $\sin^2(x)$. As before, this means there is a bijection.

\textbf{4.} We can map $0 \in [0, 1]$ to $\frac{1}{3} \in (0, 1)$, $\frac{1}{3}$ to $\frac{1}{4}$, etc. In general, let $\frac{1}{n}$ map to $\frac{1}{n + 1}$. Also, let $1 \in [0, 1]$ map to $\frac{2}{3} \in (0, 1)$, and let $\frac{n - 1}{n}$ map to $\frac{n}{n + 1}$. This goes on infinitely, leaving many numbers behind, since $(0, 1)$ is uncountable. So we can map all the remaining numbers to themselves.
\newline

\textbf{5.} Since $|S| = |T|$, there exists a bijection $f: S \rightarrow T$. Now, define a function $F: P(S) \rightarrow P(T)$ where if $A \subseteq S$ then $F(A) = \{f(x): x \in A\}$. We need to prove that this is bijective.
\newline

Suppose $A$ and $B$ are two distinct subsets of $S$. We need to show that $F(A) \neq F(B)$. For contradiction, assume that $F(A) = F(B)$. Then, by the definition of $F$, for every $x \in A$, there exists $y \in B$ such that $f(x) = f(y)$. Since $f$ is a bijection, this implies that $x = y$, which means $A = B$, which is a contradiction. Therefore, it is injective.
\newline

Now, consider a subset $Y \subseteq T$. We need to show that there exists $X \subseteq S$ such that $F(X) = Y$. Let's say $X = \{x \in S: f(x) \in Y\}$. Since $f$ is a bijection, $X \subseteq S$ and there exists $x \in S$ such that $f(x) = y$. This shows that $x \in X$ and therefore, $F(X) = Y$. So $F$ is surjective.
\newline

To show that $|P(S)| \le |P(T)|$, we need to establish an injection from $P(S)$ to $P(T)$. Since $|S| \le |T|$, there is an injection $f: S \rightarrow T$. As before, let $F(A) = {f(x): x \in A}$ where $A \subseteq S$. We need to prove that $F(A) \neq F(B)$ for two distinct subsets $A, B$ of $S$. For contradiction, assume $F(A) = F(B)$. That means, for every $x \in A$, there exists $y \in B$ such that $f(x) = f(y)$. Since $f$ is injective, $x = y$, so $A = B$ which is a contradiction. Therefore, $F$ is injective.
\newline

\textbf{6.} For any subset $A \subseteq \mathbb{N}$, consider a function $f_A(n) = 1$ if $n \in A$ and $f_A(n) = 0$ if $n \not \in A$. By this function, each subset $A$ of $\mathbb{N}$ corresponds to a binary sequence. Putting the sequence after a decimal point, we get a binary decimal between $0$ and $1$. This is an injection from the set of subsets of $\mathbb{N}$ to $[0, 1]$. It is not a bijection because of numbers like $0.111...$.
\newline

We can establish an injection from $[0, 1]$ to the set of subsets of $\mathbb{N}$ by converting it to binary, turning the numbers after the decimal point into a binary string, and reversing $f_A(n)$. Therefore, by CSB, $|P(\mathbb{N})| = c$.
\newline

\textbf{7.} Given a real number in $[0, 1]$, we can create two real numbers in $[0, 1]$ by alternating its digits and appending each of them to two separate numbers. For example, we can turn $0.49283927...$ into $0.4232...$ and $0.9897...$. We can easily reverse this process by interleaving the digits to create one number. Therefore, this is a bijection.
\newline

\textbf{8.} Let's say $A = \{f: \mathbb{Z} \rightarrow \mathbb{Z}\}$ and $B = \{f: \mathbb{Z} \rightarrow \{0, 1\}\}$. We know that $|B| = c$. We can map elements in set $B$ to themselves in set $A$, so this is an injection. Therefore, $|A| \ge |B| = c$. I couldn't figure out how to create an injection from $A$ to $B$ to then apply CSB.
\newline

\textbf{9.} The first step of constructing the Cantor set involves removing $\frac{1}{3}$ of the interval $[0, 1]$. Then, we remove $1/3$ of the remaining two intervals of length $1/3$. This gives us a length of $2/9$ that is removed. Continuing this pattern, we get a geometric sequence: $\frac{1}{3}, \frac{2}{9}, \frac{4}{27}, \ldots$. Summing them, we get $\frac{1/3}{1-\frac{2}{3}} = 1$. Therefore, a length of $1$ is removed from the interval $[0, 1]$, which leaves no length remaining. Since any nonempty open interval would have some nonzero length, it cannot exist in the Cantor set.
\newline

\textbf{10a.} If we place a figure-8 inside a bigger one, then the end-to-end length of the smaller figure-8 must be less than the diameter of the circle in the bigger figure-8 that it sits inside. That means the radii of the circles that are part of the smaller figure-8 must be less than the radius of one of the circles of the bigger figure-8 by some positive number. This means there cannot be an uncountable number of them, since they can be easily numbered outermost to innermost without skipping any.
\newline

If we place the figure-8's outside each other, there is still a nonzero distance between any two of them, so it is not uncountable. Therefore, it is impossible to place uncountably many figure-8's in the plane such that no two of them touch.
\newline

\textbf{10b.} If two circles are not concentric but one is inside the other, then there is room to include another circle tangent to it inside the bigger circle that doesn't touch it. Those two circles effectively form a figure-8, and we know from the previous problem we cannot place uncountably many of them in a plane. Therefore, the same applies for non-concentric circles.

\end{document}